\documentclass[11pt,a4paper]{article}
\usepackage[utf8]{inputenc}
\usepackage[T1]{fontenc}
\usepackage{amsmath,amssymb}
\usepackage{graphicx}
\usepackage{booktabs}
\usepackage{hyperref}
\usepackage[margin=1in]{geometry}
\usepackage{natbib}
\usepackage{xcolor}

\title{Emergent Abilities in Large Language Models: Mirage or Real?\\
\large A Re-Analysis of Published Benchmark Data}

\author{Yun Du\textsuperscript{1} \and Lina Ji\textsuperscript{1} \and Claw\textsuperscript{2}\\
\textsuperscript{1}Stanford University \quad \textsuperscript{2}AI4Science Catalyst Institute}

\date{March 2026}

\begin{document}
\maketitle

\begin{abstract}
We re-analyze published benchmark data from BIG-Bench (8 tasks, 3 model families) and MMLU (13 models, 5 families) to test the claim by \citet{schaeffer2023} that emergent abilities in large language models are artifacts of discontinuous evaluation metrics. By applying both discontinuous (exact string match) and continuous (partial credit) metrics to the same published performance data, we quantify the \emph{Metric Sensitivity Index} (MSI) for each task. We find that 7 of 8 analyzed tasks show MSI $> 2$, indicating that the apparent phase transition is primarily driven by the choice of metric. A synthetic demonstration confirms that linear per-token improvement creates dramatic apparent emergence under exact-match scoring. MMLU accuracy, a more continuous metric, scales smoothly across model families. The lone MSI $\leq 2$ case is a single-token multiple-choice task, so its MSI is definitional rather than diagnostic.
\end{abstract}

\section{Introduction}

\citet{wei2022} documented 137 tasks where large language model (LLM) performance appeared to exhibit ``emergent abilities''---capabilities absent in smaller models that appear sharply at a critical scale. This observation suggested that scaling alone could produce qualitatively new capabilities, with profound implications for AI safety and development.

\citet{schaeffer2023} challenged this interpretation, arguing that the appearance of emergence is primarily an artifact of evaluation metrics. They demonstrated that over 92\% of claimed emergent abilities are measured by just two metrics---Exact String Match and Multiple Choice Grade---both of which are discontinuous. When continuous metrics such as Token Edit Distance are applied to the \emph{same model outputs}, performance improves smoothly and predictably with scale.

We provide an independent re-analysis of this claim using hardcoded published benchmark data, requiring no model inference or API access. Our analysis introduces the Metric Sensitivity Index (MSI), a quantitative measure of how much apparent nonlinearity is attributable to metric choice versus genuine capability transitions.

\section{Methods}

\subsection{Data}

We hardcode published performance data from:
\begin{itemize}
    \item \textbf{BIG-Bench}: 8 tasks (2-digit multiplication, 4-digit addition, IPA transliteration, word unscrambling, Persian QA, sports understanding, modified arithmetic, word sorting) across GPT-3, InstructGPT, LaMDA, and PaLM model families \citep{srivastava2023,wei2022}.
    \item \textbf{MMLU}: 13 models from 5 families (GPT-3, PaLM, LLaMA, Chinchilla, Gopher) \citep{hendrycks2021,touvron2023,chowdhery2022,hoffmann2022}.
\end{itemize}

All accuracy values are in $[0, 1]$ and parameter counts in billions (B).

\subsection{Metric Transformation}

Following \citet{schaeffer2023}, we model the relationship between per-token accuracy $p$ and exact-match accuracy as:
\begin{equation}
    \text{ExactMatch} = p^n
\end{equation}
where $n$ is the number of tokens in the answer. This assumes token-level independence. The inverse yields the inferred per-token accuracy:
\begin{equation}
    p = \text{ExactMatch}^{1/n}
\end{equation}

We define three continuous metrics:
\begin{itemize}
    \item \textbf{Partial Credit}: $\text{PC} = p$ (fraction of tokens correct)
    \item \textbf{Token Edit Distance}: $\text{TED} = n(1-p)$ (expected errors)
    \item \textbf{Brier Score}: $\text{BS} = (p - y)^2$ for classification tasks
\end{itemize}

\subsection{Nonlinearity Detection}

For each task, we fit both linear and logistic sigmoid models to performance vs.\ $\log_{10}(\text{parameters})$ under both discontinuous and continuous metrics. We compare fits using $R^2$ and define the \emph{Metric Sensitivity Index}:
\begin{equation}
    \text{MSI} = \frac{(R^2_{\text{sig}} - R^2_{\text{lin}})_{\text{disc}}}{(R^2_{\text{sig}} - R^2_{\text{lin}})_{\text{cont}}}
\end{equation}

High MSI ($> 2$) indicates the nonlinearity is primarily a metric artifact. MSI $\leq 2$ suggests potentially genuine nonlinear scaling.

\subsection{Synthetic Demonstration}

We generate synthetic data where per-token accuracy improves linearly with $\log_{10}(\text{parameters})$ from $p=0.3$ to $p=0.95$ across 20 simulated model sizes. We then apply both exact-match ($p^5$) and partial-credit scoring to demonstrate how the nonlinear mapping creates apparent emergence.

\section{Results}

\subsection{Metric Sensitivity Index}

Of 8 BIG-Bench tasks analyzed, \textbf{7 show MSI $> 2$}, indicating that the apparent phase transition is predominantly a metric artifact (Table~\ref{tab:msi}). The remaining task, sports understanding, has a single-token output under multiple-choice accuracy, so exact match and per-token accuracy are identical by construction and MSI is not informative for that case.

\begin{table}[h]
\centering
\caption{Metric Sensitivity Index for BIG-Bench tasks. MSI $> 2$ indicates likely metric artifact.}
\label{tab:msi}
\begin{tabular}{lcccc}
\toprule
\textbf{Task} & \textbf{MSI} & \textbf{$R^2_{\text{sig}}$ (EM)} & \textbf{$R^2_{\text{lin}}$ (PC)} & \textbf{Verdict} \\
\midrule
2-Digit Multiplication & 27.06 & 0.864 & 0.549 & Artifact \\
4-Digit Addition & 55.17 & 0.753 & 0.595 & Artifact \\
IPA Transliterate & 4.78 & 0.986 & 0.979 & Artifact \\
Word Unscramble & 1401.51 & 0.978 & 0.846 & Artifact \\
Modified Arithmetic & 7.82 & 0.978 & 0.868 & Artifact \\
Word Sorting & 6.57 & 0.979 & 0.969 & Artifact \\
Persian QA & 7.69 & 0.993 & 0.979 & Artifact \\
Sports Understanding & 1.00 & 0.890 & 0.890 & N/A (single token) \\
\bottomrule
\end{tabular}
\end{table}

\subsection{Synthetic Demonstration}

The synthetic demonstration confirms the core mechanism. With 5-token answers and linearly improving per-token accuracy ($p = 0.30 \to 0.95$):
\begin{itemize}
    \item Partial credit improves smoothly from 0.30 to 0.95
    \item Exact match ($p^5$) shows near-zero performance below $p \approx 0.7$, then rises sharply
    \item The same underlying improvement appears as smooth scaling or dramatic emergence depending solely on the metric
\end{itemize}

\subsection{MMLU Scaling}

MMLU accuracy (5-shot, multiple-choice) scales relatively smoothly across all model families. Linear $R^2$ values are high for within-family scaling (GPT-3: $R^2 > 0.9$; LLaMA: $R^2 > 0.99$), consistent with the absence of phase transitions when using a more continuous evaluation metric.

\section{Discussion}

Our re-analysis provides evidence consistent with \citet{schaeffer2023}: most apparent emergent abilities are metric artifacts caused by the nonlinear mapping $p \to p^n$ inherent in exact-match scoring. The Metric Sensitivity Index provides a quantitative framework for distinguishing genuine capability transitions from metric artifacts.

However, we note several important caveats:
\begin{enumerate}
    \item \textbf{Sparse data}: With only 3--14 model sizes per task, curve-fitting comparisons have limited statistical power.
    \item \textbf{Token independence}: Our per-token accuracy inference assumes independence, which may not hold for reasoning-intensive tasks.
    \item \textbf{Aggregated scores}: We use published accuracy values, not raw model outputs, preventing direct verification of the per-token distribution.
    \item \textbf{Hardcoded data}: All data is transcribed from published figures and tables, introducing potential transcription error.
\end{enumerate}

Sports understanding should not be treated as counterevidence to the metric-artifact hypothesis: with a single-token output, exact match and partial credit are the same metric. This means our current MSI analysis cannot adjudicate whether that task contains genuine nonlinearity, and future work should use task-level metrics that remain informative when $n=1$.

\section{Conclusion}

We confirm the central finding of \citet{schaeffer2023} using an independent re-analysis: 7 of 8 BIG-Bench tasks show high Metric Sensitivity Index, indicating that apparent emergence is primarily a metric artifact. The Metric Sensitivity Index provides a useful quantitative tool for future studies of scaling behavior, and our fully reproducible analysis requires no model inference.

\bibliographystyle{plainnat}
\begin{thebibliography}{99}

\bibitem[Chowdhery et al., 2022]{chowdhery2022}
Chowdhery, A., et al. (2022).
\newblock PaLM: Scaling Language Modeling with Pathways.
\newblock \emph{arXiv:2204.02311}.

\bibitem[Hendrycks et al., 2021]{hendrycks2021}
Hendrycks, D., et al. (2021).
\newblock Measuring Massive Multitask Language Understanding.
\newblock \emph{ICLR 2021}. arXiv:2009.03300.

\bibitem[Hoffmann et al., 2022]{hoffmann2022}
Hoffmann, J., et al. (2022).
\newblock Training Compute-Optimal Large Language Models.
\newblock \emph{arXiv:2203.15556}.

\bibitem[Schaeffer et al., 2023]{schaeffer2023}
Schaeffer, R., Miranda, B., \& Koyejo, S. (2023).
\newblock Are Emergent Abilities of Large Language Models a Mirage?
\newblock \emph{NeurIPS 2023}. arXiv:2304.15004.

\bibitem[Srivastava et al., 2023]{srivastava2023}
Srivastava, A., et al. (2023).
\newblock Beyond the Imitation Game: Quantifying and Extrapolating the Capabilities of Language Models.
\newblock \emph{arXiv:2206.04615}.

\bibitem[Touvron et al., 2023]{touvron2023}
Touvron, H., et al. (2023).
\newblock LLaMA: Open and Efficient Foundation Language Models.
\newblock \emph{arXiv:2302.13971}.

\bibitem[Wei et al., 2022]{wei2022}
Wei, J., et al. (2022).
\newblock Emergent Abilities of Large Language Models.
\newblock \emph{arXiv:2206.07682}.

\end{thebibliography}

\end{document}
